\section{Orchestration}
\label{sec:orchestration}

\subsection{The variation step as a skill}

The \cmd{avo-vary} skill tells the agent: ``You are the variation operator \ldots\ You do
\emph{not} decide what gets kept --- \avo{commit} does.'' It prescribes one step in a fixed order:
\begin{lstlisting}
avo mem prime                 # 1. what past turns learned, incl. dead ends
avo best; avo lineage         # 2. the version to beat, and the whole population
avo know query "<the idea>"   # 3. K: docs and prior lineage entries
# 4. make ONE coherent change
avo score --json              # 5. measure it
avo commit --why "<what changed and why it should help>"
\end{lstlisting}
Steps 1--3 are the agent's reading of $\Pt$ and $\K$ in \cref{eq:intro-agent}. The rule ``one idea
per candidate'' makes a refusal attributable to one change. Inside this step the agent is free: it
can profile, read code, search the web, or call \avo{fan} (\cref{sec:fan}). This is the paper's
edit--evaluate--diagnose cycle, written as instructions rather than code.

\subsection{\texorpdfstring{\avo{run}}{avo run}: the continuous loop}
\label{sec:run}

\avo{run} automates the cycle of \cref{fig:loop}. Its central choice differs from the paper:
\emph{each turn is a fresh agent process}. The paper's agent keeps one conversation for seven days.
In avocode, ``every iteration is a fresh process, so the turn prompt is the only continuity''. All
other state lives in the repository: the lineage in git, the trajectory in the attempt log, and
memory in beads or JSONL. \Cref{alg:run} gives the loop.

\begin{algorithm}[t]
\caption{\avo{run}: the outer loop}
\label{alg:run}
\begin{algorithmic}[1]
\Require repo with seed commit $x_0$; prompt $\pi$; agent $\alpha$; budget $T$; turn timeout $\tau$;
  thresholds $N$ (stall), $R$ (thrash)
\State \Call{CheckGuards}{$\pi$} \Comment{depth $<$ cap, $\mathrm{sha}(\pi) \notin$ chain}
\State $d \gets \bot$;\ $\mathit{prev} \gets \bot$;\ $\mathit{idle} \gets 0$
\For{$t = 1, \ldots, T$}
  \If{\file{.avo/STOP} exists} \Return \textsc{stop-file} \EndIf
  \State $h_0 \gets \texttt{HEAD}$
  \State $p \gets \Call{TurnPrompt}{\pi, t, \mathit{prev}, d}$ \Comment{$\pi$ verbatim at $t=1$}
  \State $(\mathit{ok}, \mathit{summary}) \gets \Call{DriveAgent}{\alpha, p, \tau}$
    \Comment{fresh process: $\Vary(\Pt)$}
  \If{spawn failed} \Return \textsc{agent-unavailable} \EndIf
  \State $\mathit{dec} \gets \Call{Commit}{\mathit{summary}}$ \Comment{\cref{alg:commit}; summary becomes \emph{why}}
  \State $h_1 \gets \texttt{HEAD}$;\ credit versions in $h_0..h_1$ \Comment{agent may commit itself}
  \State $\mathit{idle} \gets (\mathit{dec} = \textsc{noop} \wedge h_1 = h_0)\ ?\ \mathit{idle}+1 : 0$
  \State $(\mathit{sig}, d) \gets \Call{Supervise}{N, R}$ \Comment{\cref{alg:supervise}}
  \If{$\mathit{sig} \neq \emptyset$} record an \cmd{intervention} memory \EndIf
  \State write manifest \Comment{after every iteration}
  \State $\mathit{prev} \gets (\mathit{dec}, \mathit{sig})$
  \If{$\mathit{idle} \geq 3$} \Return \textsc{no-progress} \EndIf
\EndFor
\State \Return \textsc{max-iters}
\end{algorithmic}
\end{algorithm}

\begin{algorithm}[t]
\caption{\textsc{Commit}($\mathit{why}$): the only writer of $\Pt$}
\label{alg:commit}
\begin{algorithmic}[1]
\If{tree is clean, ignoring harness paths} \Return \textsc{noop} \EndIf
\State $a \gets \f(\text{working tree})$;\ append $a$ to the attempt log
\State $b \gets$ newest version in $\Pt$ \Comment{best, by \cref{prop:monotone}}
\If{$a$ has harness errors} \Return \textsc{error} \EndIf
\If{$\neg\pass(a)$ \textbf{or} ($b \neq \bot$ \textbf{and not} \Call{Accept}{$a, b$})} \Comment{\cref{def:commit}}
  \State record dead end, key $= \mathrm{sha256}(\mathit{reason}\,\|\,\mathit{vector})$
  \State \Return \textsc{refused}
\EndIf
\State commit with trailers; add git note; write \file{lineage/vNNN.md}; record version memory
\State \Return \textsc{committed}
\end{algorithmic}
\end{algorithm}

\paragraph{The turn prompt.} The first turn gets the operator's prompt verbatim. Later turns get the
prompt plus a footer: the turn number, a one-line outcome of the previous turn (for example
\emph{``committed v3 at 0.408 ms''} or \emph{``refused at 0.42 ms --- \ldots''}), the number of
attempts since the best version, a reminder that the agent is a fresh process that should read
\avo{lineage}, \avo{best} and \avo{mem prime}, and the supervisor's directive if one fired. The
directive is \emph{appended} to the task, never substituted for it.

\paragraph{Committing.} The agent's final message, capped at 2{,}000 characters, becomes the
\emph{why} of the commit. The skill also allows the agent to call \avo{commit} itself during the
turn. The harness therefore reads the trailers in $h_0..h_1$ and credits those versions too. Before
this fix, a run that grew four versions reported that it had committed none.

\paragraph{Termination.} The loop stops on the iteration budget, on a \file{.avo/STOP} file, after
three consecutive idle turns, or when the agent binary cannot start. The stop file needs no process
to find, and the agent can create it when the task is done. The idle rule exists because ``an
unchanged tree records no attempt, so the stall detector cannot see it''. A timed-out turn does not
stop the loop; its process group is killed and the next turn starts. The supervisor never stops the
loop. It only steers.

\paragraph{Crash safety.} The run manifest is the same object as the \cmd{--json} report, and it is
rewritten after every iteration: ``a days-long loop gets killed eventually, and per-iteration versus
at-exit is a recoverable record versus nothing''. There is no \cmd{--resume}. Recovery is structural:
each version is committed before the next turn, so a new \avo{run} with the same prompt continues
from \cmd{HEAD}. The second benchmark run in \cref{sec:eval} did exactly this.

\paragraph{Budgets.} \avo{run} has an iteration budget and a per-turn timeout (default 900\,s). It
has no wall-clock budget and no cost budget. Cost is recorded but not enforced.

\subsection{\texorpdfstring{\avo{supervise}}{avo supervise}: stagnation detection}
\label{sec:supervise}

The paper's supervisor is an agent. avocode's supervisor is a pure function of what the harness
already records, plus a text template. It makes no LLM call. The design reason is in the source:
\begin{quote}\small
``The output is a directive that \textbf{cites}. `Try something else' is worthless --- the agent
already believes it is.''
\end{quote}
It detects two signals. Let $A = (a_1, \ldots, a_m)$ be the last 1{,}000 attempts, newest last, and
$b$ the best version.

\begin{definition}[Stall]
Let $s = |\{a \in A : \mathrm{head}(a) = \mathrm{sha}(b) \vee \mathrm{ts}(a) \geq \mathrm{date}(b) +
1\,\mathrm{s}\}|$, or $s = m$ if no version exists. Stall fires when $s \geq N$ and $s > 0$, with
$N = 5$ by default.
\end{definition}
\noindent Attempts are matched to $b$ by git \cmd{HEAD} first, not by clock: the attempt that
\emph{became} $b$ was scored before the commit and carries the parent's SHA, and git truncates
commit dates to the second. An attempt that passed $\f$ but did not beat $b$ still counts.

\begin{definition}[Thrash]
Let $\mathrm{sig}(a)$ be the failure signature of $a$, or $\bot$ for a passing attempt. Let $r$ be the
length of the longest suffix of $A$ whose signatures all equal $\mathrm{sig}(a_m) \neq \bot$. Thrash
fires when $r \geq R$, with $R = 3$ by default.
\end{definition}
\noindent The signature is the first harness error, else the scorer's first log line, normalized so
that noise does not split one failure into many: lowercase; temporary paths, hex SHAs and numbers
replaced by placeholders; whitespace folded; truncated to 200 characters. A pass breaks the streak.
Thrash ``usually means the diagnosis is wrong rather than the edit''.

There is no test for a score plateau, for repeated ideas, or for similarity between a candidate and
prior versions. The harness ranks measurements, not ideas. Novelty is left to the agent, which the
directive gives the evidence to judge.

\paragraph{The directive.} Only when a signal fires does the supervisor read memory and $\K$, so the
common case stays cheap. The directive cites:
\begin{itemize}[leftmargin=*,itemsep=1pt]
  \item up to three recent versions, with score, rationale and lineage file (``a variation that only
    re-derives one of these is not progress'');
  \item up to five memories, dead ends before insights;
  \item up to three documents in $\K$ that no version or memory has mentioned.
\end{itemize}
A document $d$ with significant title terms $T(d)$ counts as \emph{unexplored} against the corpus
$C$ of all version rationales and memories when
\begin{equation}
  2\,\big|\{u \in T(d) : u \in C\}\big| < |T(d)|.
\end{equation}
The bar is low on purpose: a document wrongly marked explored costs one citation, while one wrongly
marked unexplored ``sends the agent back down a road it already travelled''. Interventions are never
cited. Citing them would nest old directives inside new ones, and would put every recommended
document name into the corpus, so ``it would erase its own best advice''. \Cref{fig:directive} shows
a real directive from a test fixture.

\begin{figure}[t]
\begin{lstlisting}
STEERING (avo supervise)

- stall: 5 attempt(s) since v1 (16 bytes) with no committed improvement; the threshold is 5

Already committed (a variation that only re-derives one of these is not progress):
  - v1 - 16 bytes - shrink the kernel by removing five redundant bytes (lineage/v001.md)

In K and never mentioned by any version - unexplored:
  - knowledge/loop-tiling.md - Loop tiling and cache blocking
  - knowledge/register-pressure.md - Register pressure on Blackwell

Do this next:
  1. Name the direction you are about to take, and say how it differs from every line
     above. If you cannot, you do not have a new idea yet.
  2. Change one thing, then `avo score` before you change a second.
  3. `avo commit --why "<the direction>"` when it measures better. A refusal is a
     measurement: record it with `avo mem add`, do not retry it.
  4. If every direction you can name is already above, do not guess: `avo fan --n 4`
     explores four at once and lets f choose.
\end{lstlisting}
\caption{A stall directive (abridged), generated in a fixture repository with one version, five
attempts that did not beat it, and two documents in $\K$. When thrash fires, step 1 instead says
that ``the same error $N$ times means the \emph{diagnosis} is wrong, not the edit''.}
\label{fig:directive}
\end{figure}

\begin{algorithm}[t]
\caption{\textsc{Supervise}($N, R$)}
\label{alg:supervise}
\begin{algorithmic}[1]
\State $A \gets$ last 1000 attempts;\ $b \gets$ newest version
\State $s \gets |\{a \in A : \mathrm{head}(a) = \mathrm{sha}(b) \vee \mathrm{ts}(a) \geq \mathrm{date}(b) + 1\mathrm{s}\}|$
\State $r \gets$ longest suffix of $A$ with one non-null signature
\State $\mathit{sig} \gets \{\textsc{stall} \mid s \geq N \wedge s > 0\} \cup \{\textsc{thrash} \mid r \geq R\}$
\If{$\mathit{sig} = \emptyset$} \Return $(\emptyset, \bot)$ \EndIf
\State $V \gets$ newest 3 versions;\ $F \gets$ up to 5 memories (dead ends first, no interventions)
\State $U \gets$ up to 3 documents in $\K$ that are unexplored
\State \Return $(\mathit{sig}, \Call{Render}{\mathit{sig}, V, F, U})$ \Comment{$\leq 4000$ characters}
\end{algorithmic}
\end{algorithm}

\paragraph{Two delivery paths.} In \avo{run}, the directive enters the \emph{next} turn's prompt and
is recorded as an \cmd{intervention} memory. In the pi agent, an extension runs the same detector
after every \cmd{avo\_score} or \cmd{avo\_commit} tool result and injects the directive into the
\emph{current} session as a steering message. It sends one directive per episode: an episode key
built from the best version's SHA and the start of the streak makes ``5 since best'' and ``6 since
best'' the same episode. Its state is rebuilt from the session branch, so a branched session that
never saw a directive gets one. Both paths count the attempt log, not the session: ``two drifting
counters would steer one stall twice''.

\subsection{\texorpdfstring{\avo{fan}}{avo fan}: parallel probes}
\label{sec:fan}

\avo{fan} explores $N$ directions at once. Each probe gets its own detached git worktree at
\cmd{HEAD} and its own headless agent process. The probes run on a small model (set by
\cmd{AVO\_PROBE\_MODEL}), because ``exploration is a small-model job, exploitation is a big-model
job''. \avo{run} uses the agent's default model. All probes get the \emph{same} prompt; diversity
comes from sampling. The skill therefore asks for a prompt that names the problem, not the edit:
``delete the bounds check'' gives $N$ identical diffs.

Each probe is scored in its worktree without writing to the root attempt log, so probes never
trigger the supervisor. A probe is scored even if its agent failed, since ``a half-finished edit
passing $\f$ is a real result''. Worktrees with no changes are removed. The manifest is written after
every probe, so \cmd{--resume} reruns only unfinished probes after a crash.

\paragraph{Selection.} Let $\nu(x) = \mathrm{primary}(x)$ if higher is better and $-\mathrm{primary}(x)$
otherwise. The harness computes
\begin{equation}
  \hat{\imath} = \operatorname*{arg\,max}_{i\,:\,\pass(x_i)} \nu(x_i),
  \qquad \hat{\imath} = \bot \text{ if no probe passes},
\end{equation}
with ties going to the smallest index. This is ``a hint, never a decision''. Promotion is manual:
\avo{fan --promote} $j$ saves and applies probe $j$'s diff to the root tree and stops. The candidate
must then go through \avo{score} and \avo{commit}, which re-measures it in the real tree:
\begin{equation}
  x_{t+1} =
  \begin{cases}
    \mathrm{promote}(x_j) & \text{if the commit rule accepts it against } x_t,\\
    x_t & \text{otherwise.}
  \end{cases}
\end{equation}
The skill allows at most one promotion per step, because ``two probes' diffs stacked together is a
candidate neither of them measured''.

\paragraph{Recursion guards.} A probe is an agent with the same skills, so it can call \avo{fan}
itself, and a turn of \avo{run} can call \avo{run}. Four guards bound this, carried in environment
variables, ``the only channel surviving \cmd{spawn} into an arbitrary agent binary'':
\begin{enumerate}[leftmargin=*,itemsep=1pt]
  \item \textbf{Depth}: refuse when the level reaches the cap (default 3); a child gets level $+1$.
  \item \textbf{Cycle}: refuse when the prompt's hash is already in the chain of ancestors.
  \item \textbf{Concurrency}: at most $\min(8, \mathit{cpus} - 2)$ probes at once.
  \item \textbf{Timeout}: 900\,s per probe, after which the whole process group is killed.
\end{enumerate}
A guard refusal is exit 1, not 2: ``a guard is a refusal, not a harness failure''. avocode uses OS
processes and worktrees, not in-process sub-agents, because the design notes judge that sub-agent
extensions ``do not work well enough''.

\subsection{Agent adapters}
\label{sec:adapters}

One function, \cmd{driveAgent}, runs one headless turn for both \avo{run} and \avo{fan}.
\Cref{tab:adapters} lists the templates. Every flag was read off the real binaries, because ``a
wrong flag fails silently''. The prompt follows \cmd{--}, so a prompt that starts with a dash is not
read as a flag. Token counts are normalized to four disjoint fields: uncached input, output, cache
read and cache write. Codex reports cached input as a subset of input, so it is subtracted.

\begin{table}[t]
\centering\small
\caption{Headless agent templates. A custom template can be set in \file{.avo/config.json}.}
\label{tab:adapters}
\begin{tabularx}{\linewidth}{@{}l X X@{}}
\toprule
Agent & Command & Skill discovery \\
\midrule
pi \citep{zechner2025pi} & \cmd{pi --mode json --print --approve -- <prompt>} & native \file{.agents/skills/}; six native tools; in-session supervisor \\
Claude Code \citep{anthropic2025claudecode} & \cmd{claude --print --output-format stream-json --verbose --permission-mode bypassPermissions -- <prompt>} & \file{.claude/skills} $\to$ \file{.agents/skills} symlink \\
Codex \citep{openai2025codex} & \cmd{codex exec --json --skip-git-repo-check --sandbox workspace-write -- <prompt>} & skills index in \file{AGENTS.md} \\
\bottomrule
\end{tabularx}
\end{table}

\avo{install} wires the skills into a target repository by \emph{linking}, never copying: ``a copy is
a fork that stops receiving fixes''. It validates each skill against the Agent Skills
specification first, because pi is lenient, Claude Code is strict, and a skill must load everywhere.
It also writes a managed block into \file{AGENTS.md} with the rules and a skills index, and leaves
text outside the markers alone.

For pi, an extension adds six native tools (\cmd{avo\_score}, \cmd{avo\_commit},
\cmd{avo\_lineage}, \cmd{avo\_know\_query}, \cmd{avo\_know\_add}, \cmd{avo\_fan}). They are thin
wrappers over the same functions as the CLI. A refusal is a normal result, not an error. No tool
takes a working directory, because ``an agent that can retarget the repo can write a version into a
repo nobody is watching''.
